% LineContainer (CHT)

Maintains a dynamic set of lines $y = kx + m$ to query the maximum or minimum value at any $x$ in $O(\log N)$ using a variant of \lstinline|std::set|.

\begin{lstlisting}
struct Line {
    mutable long long k, m, p;
    bool operator<(const Line& o) const { return k < o.k; }
    bool operator<(long long x) const { return p < x; }
};

struct LineContainer : set<Line, less<>> {
    static const long long INF = LLONG_MAX;
    
    long long div(long long a, long long b) {
        return a / b - ((a ^ b) < 0 && a % b);
    }
    
    bool isect(iterator x, iterator y) {
        if (y == end()) { x->p = INF; return false; }
        if (x->k == y->k) x->p = x->m > y->m ? INF : -INF;
        else x->p = div(y->m - x->m, x->k - y->k);
        return x->p >= y->p;
    }
    
    void add(long long k, long long m) {
        auto z = insert({k, m, 0}), y = z++, x = y;
        while (isect(y, z)) z = erase(z);
        if (x != begin() && isect(--x, y)) { isect(x, y = erase(y)); }
        while ((y = x) != begin() && (--x)->p >= y->p) isect(x, erase(y));
    }
    
    long long query(long long x) {
        assert(!empty());
        auto l = *lower_bound(x);
        return l.k * x + l.m;
    }
};
\end{lstlisting}

This structure natively computes the \textbf{MAXIMUM} value. 
\begin{itemize}
    \item \textbf{To find Maximum:} Call \lstinline|add(k, m)| and query directly via \lstinline|query(x)|.
    \item \textbf{To find Minimum:} Insert lines as \lstinline|add(-k, -m)|. When querying, negate the output: \lstinline|-query(x)|.
\end{itemize}
